\documentclass[11pt,a4paper]{article}

% =========================
% Encoding & Fonts (pdflatex-safe)
% =========================
\usepackage[utf8]{inputenc}
\usepackage[T1]{fontenc}
\usepackage{lmodern}

% =========================
% Layout
% =========================
\usepackage[a4paper,margin=1.5cm]{geometry}
\setlength{\parindent}{0pt}
\setlength{\parskip}{6pt}

% =========================
% Packages
% =========================
\usepackage{amsmath,amssymb}
\usepackage{enumitem}
\usepackage{hyperref}
\usepackage{xcolor}
\usepackage{listings}
\usepackage{titlesec}
\usepackage{fancyhdr}
\usepackage{tabularx}
\usepackage{array}
\usepackage{graphicx}

% =========================
% Header / Footer
% =========================
\pagestyle{fancy}
\fancyhf{}
\lhead{\textbf{Lebanese University -- Faculty of Engineering}}
\rhead{\textbf{Semester VIII}}
\cfoot{\thepage}

% =========================
% Section formatting
% =========================
\titleformat{\section}{\Large\bfseries}{\thesection.}{0.5em}{}
\titleformat{\subsection}{\large\bfseries}{\thesubsection.}{0.5em}{}
\titleformat{\subsubsection}{\normalsize\bfseries}{\thesubsubsection.}{0.5em}{}

% =========================
% Code style
% =========================
\lstset{
	basicstyle=\ttfamily\small,
	frame=single,
	breaklines=true,
	columns=fullflexible,
	keywordstyle=\color{blue},
	commentstyle=\color{gray},
	stringstyle=\color{red},
	showstringspaces=false
}
% ---------------- Colors ----------------
\definecolor{primary}{HTML}{1F4E79}
\definecolor{secondary}{HTML}{333333}
% =========================
% Helpers
% =========================
\newcommand{\courseTitle}{Concurrency \& Distributed Systems}
\newcommand{\courseSubtitle}{Diagnosing Broken Systems: From Threads to Clusters}
\newcommand{\primaryLang}{Java}
\newcommand{\counterWeek}{Python (with brief Go/Erlang contrast)}
\newcommand{\level}{Senior Undergraduate / Graduate}
\newcommand{\format}{Theory + Failure-Driven Labs + Evolving Project + Presentations}
\newcommand{\instructor}{Dr. Mohamad Aoude}
\usepackage{fancyhdr}
\usepackage{lastpage}
\usepackage{graphicx}
\usepackage{eso-pic}
\usepackage{qrcode}

% -------------------------------
% Header & Footer
% -------------------------------
\pagestyle{fancy}
\fancyhf{}

\fancyhead[L]{\textbf{ULFG --- SEM VIII}}
\fancyhead[R]{\textbf{Dr M AOUDE}}

\fancyfoot[L]{\textcopyright\ \the\year\ M AOUDE -- All Rights Reserved}
\fancyfoot[C]{Page \thepage\ /\ \pageref{LastPage}}
\fancyfoot[R]{\today}

\renewcommand{\headrulewidth}{0.4pt}
\renewcommand{\footrulewidth}{0.4pt}
\begin{document}
	% ---------------- Title Page ----------------
\begin{titlepage}
	\centering
	\vspace*{3cm}
	
	{\Huge\bfseries\color{primary}
		Concurrency, Parallelism,\\
		Distributed Systems}
	
	\vspace{1cm}
	{\Large Systems-First Detailed Course Syllabus}
	
	\vspace{2cm}
	{\large
		Lebanese University\\
		Faculty of Engineering\\
		Semester VIII\\
		Academic Year 2025--2026
	}
	
	\vfill
	{\large
		Instructor:\\
		\textbf{Dr.\ Mohamad Aoude}
	}
	
	\vspace{1cm}
\end{titlepage}	
	% =========================
	% Cover Block
	% =========================
	\begin{center}
		{\Large \textbf{Lebanese University}}\\[1mm]
		{\large \textbf{Faculty of Engineering}}\\[4mm]
		
		{\LARGE \textbf{\courseTitle}}\\[1mm]
		{\large \textbf{\courseSubtitle}}\\[4mm]
		
		\begin{tabular}{>{\bfseries}rl}
			Primary Language: & \primaryLang \\
			Counter-Example Week: & \counterWeek \\
			Level: & \level \\
			Format: & \format \\
			Semester: & VIII \\
			Instructor: & \instructor \\
		\end{tabular}
	\end{center}
	
	\vspace{4mm}
	\hrule
	\vspace{6mm}
	
	% ============================================================
	\section*{Course Philosophy}
	
	\begin{quote}
		\textbf{This course teaches principles, not products.}\\
		Tools will change. The problems will not.\\
		You are learning to diagnose illnesses --- not memorize prescriptions.
	\end{quote}
	
	% ============================================================
	\section*{Course Promise (Week 0)}
	
	\begin{quote}
		\textbf{By the end of this course, you will be able to diagnose and fix real concurrent system bugs professionally.}\\
		This is not theoretical knowledge --- it is a marketable engineering skill.
	\end{quote}
	
	% ============================================================
	\section*{Learning Outcomes}
	
	By the end of the course, students will be able to:
	
	\begin{itemize}[leftmargin=*]
		\item Distinguish concurrency, parallelism, and distribution
		\item Design systems that are thread-safe, visibility-correct, deadlock-free, and live
		\item Use professional debugging tools (thread dumps, profilers) to diagnose failures
		\item Test concurrent code realistically (stress, timing control, reproduction strategy)
		\item Build and evolve a server from single-threaded $\rightarrow$ pooled $\rightarrow$ async $\rightarrow$ distributed
		\item Design distributed systems under latency, partial failure, retries, and backpressure
		\item Make defensible architectural decisions (threads vs async vs Kafka vs RPC vs sockets)
	\end{itemize}
	
	% ============================================================
	\section*{Semester Roadmap --- What You Will Build}
	
	You will build \textbf{one system}, evolving it across the semester:
	
	\begin{verbatim}
		[ Week 3 ]   Single-threaded  ->  Thread-safe
		(correctness under concurrency)
		
		[ Week 5 ]   Thread pool + Backpressure
		(bounded resources, fairness)
		
		[ Week 8 ]   Async + Testing
		(non-blocking design, reproducibility)
		
		[ Week 14 ]  Distributed + Failure Recovery
		(timeouts, retries, tracing, resilience)
	\end{verbatim}
	
	Each stage \textbf{reuses the previous one} --- no throwaway work.
	
	\begin{quote}
		\textbf{Rule of the course:}\\
		\textit{You don't move forward by adding features --- you move forward by surviving failure.}
	\end{quote}
	
	% ============================================================
	\section*{Semester Project Arc: What You'll Build (From Day 1)}
	
	A single evolving system, built in stages:
	
	\begin{enumerate}[leftmargin=*]
		\item \textbf{Week 3 (Checkpoint 1):} Thread-safe server (correctness under concurrency)
		\item \textbf{Week 5 (Checkpoint 2):} Thread-pooled server (bounded resources + fairness)
		\item \textbf{Week 8 (Checkpoint 3 / Midterm):} Async server (non-blocking design + testing)
		\item \textbf{Week 14 (Checkpoint 4 / Final):} Distributed cluster (timeouts, tracing, idempotency, resilience)
	\end{enumerate}
	
	% ============================================================
	\section*{Core Framework: The Concurrency Scorecard (Used Weekly)}
	
	Every lab submission must include a completed scorecard:
	
	\begin{verbatim}
		[ ] Thread-safe? (no races)
		[ ] Visibility guaranteed? (JMM reasoning)
		[ ] Deadlock-free? (lock discipline)
		[ ] Liveness guaranteed? (no starvation/livelock)
		[ ] Bounded resources? (queues/pools/limits)
		[ ] Failure recovery path? (timeouts/retries/fallbacks)
	\end{verbatim}
	
	% ============================================================
	\section*{Explanation Quality Rubric (Used in All Reports)}
	
	\begin{itemize}[leftmargin=*]
		\item \textbf{Poor:} ``It's synchronized / it works now.''
		\item \textbf{Good:} ``Synchronization enforces mutual exclusion on this monitor, preventing interleaving that violates invariant X.''
		\item \textbf{Excellent:} ``This establishes happens-before between operations A and B, preventing reordering/visibility issues, preserving invariant X under schedule Y; tradeoff is contention cost Z.''
	\end{itemize}
	
	% ============================================================
	\section*{Policies (Because This Is an Iterative Engineering Course)}
	
	\begin{itemize}[leftmargin=*]
		\item \textbf{Late policy:} $-10\%$ per day, max 3 days late unless approved.
		\item \textbf{Revision policy:} Each checkpoint \textbf{can be revised once} after feedback, due \textbf{7 days after graded return}.
		\item \textbf{Collaboration:} Discuss ideas freely; code and reports must be original.
	\end{itemize}
	

	
	\begin{quote}
		\textbf{Professional rule:}\\
		\textit{If someone else can lock your object, your invariants are already broken.}
	\end{quote}
	
	% ============================================================
	\section*{ --- Readiness \& Alignment}
	

	\begin{itemize}[leftmargin=*]

		\item Recommended review resources:
		\begin{itemize}[leftmargin=*]
			\item Java threads + basic I/O primer (instructor-provided links/chapters)
			\item Basic CLI + Git workflow quickstart
		\end{itemize}
	\end{itemize}
	
	% ============================================================
	\section*{Phase 1 --- Shared-Memory Concurrency (Weeks 1--8)}
	
	% ------------------------------------------------------------
	\subsection*{Week 1 --- Why Concurrency Exists (Pain First)}
	
	\subsubsection*{Motivation Demo}
	Single-threaded server: one request at a time $\rightarrow$ queueing $\rightarrow$ high latency $\rightarrow$ idle CPU during I/O.
	
	\subsubsection*{Concepts}
	\begin{itemize}[leftmargin=*]
		\item Why concurrency exists
		\item Concurrency vs parallelism
		\item CPU vs cores
		\item I/O wait vs CPU work
		\item Amdahl's Law 
	\end{itemize}
	
	\subsubsection*{Lab 1.1}
	Single-threaded vs multi-threaded \textbf{blocking} server.
	
	\textbf{Deliverables:}
	\begin{itemize}[leftmargin=*]
		\item Measurements (latency/throughput)
		\item Completed scorecard
		\item Short explanation using rubric
	\end{itemize}
	
	\subsubsection*{Case Study (Mini)}
	Redis/SQLite: when single-threaded wins (simplicity + cache locality).
	
	\subsubsection*{Common Mistakes}
	\begin{itemize}[leftmargin=*]
		\item ``More threads = faster''
		\item Confusing throughput with responsiveness
	\end{itemize}
	
	% ------------------------------------------------------------
	\subsection*{Week 2 --- Threads, Lifecycle, and When NOT to Use Concurrency}
	
	\subsubsection*{Concepts}
	\begin{itemize}[leftmargin=*]
		\item \texttt{Thread} vs \texttt{Runnable}
		\item Thread states: NEW, RUNNABLE, BLOCKED, WAITING
		\item Scheduling myths
		\item When single-thread designs win (contention, predictability)
	\end{itemize}
	
	\subsubsection*{Lab 2.1}
	Thread creation patterns + join vs sleep.
	
	\textbf{Deliverables:}
	\begin{itemize}[leftmargin=*]
		\item Completed scorecard
		\item Short failure analysis: ``why sleep isn't sync''
	\end{itemize}
	
	\subsubsection*{Pitfalls Cheat Sheet (Release 1)}
	\begin{itemize}[leftmargin=*]
		\item Don't use \texttt{sleep()} for synchronization
		\item Don't assume order/fairness
	\end{itemize}
	
	% ------------------------------------------------------------
	\subsection*{Week 3 --- Synchronization \& Locks (The Awakening)}
	
	\subsubsection*{Concepts}
	\begin{itemize}[leftmargin=*]
		\item Race conditions, critical sections
		\item \texttt{synchronized}, intrinsic locks
		\item Debug tools: \texttt{jstack}, thread dumps, VisualVM (basic)
	\end{itemize}
	
	\subsubsection*{Deadlock Taxonomy}
	\begin{enumerate}[leftmargin=*]
		\item Deadly Embrace
		\item Resource Starvation
		\item Livelock
		\item Phantom Deadlock
	\end{enumerate}
	
	\subsubsection*{Lab 3.1}
	Bank account race condition $\rightarrow$ fix with correct locking.
	
	\subsubsection*{Lab 3.2 (Concrete lock-order scenario)}
	\textbf{Transfer between accounts A$\rightarrow$B and B$\rightarrow$A}
	\begin{itemize}[leftmargin=*]
		\item Create deadlock intentionally
		\item Fix using \textbf{lock ordering by account ID}
	\end{itemize}
	
	\subsubsection*{Lab 3.3 --- Deadlock + Thread Interruption (Failure Is Concurrent)}
	\textbf{Extension to Deadlock Lab:} after detecting the deadlock:
	\begin{itemize}[leftmargin=*]
		\item Interrupt one of the blocked threads using \texttt{Thread.interrupt()}
		\item Require students to:
		\begin{itemize}[leftmargin=*]
			\item Catch \texttt{InterruptedException}
			\item \textbf{Restore interrupt status} (\texttt{Thread.currentThread().interrupt()})
			\item Exit cleanly without corrupting shared state
		\end{itemize}
	\end{itemize}
	
	\textbf{Explicit Rule:}
	\begin{quote}
		Swallowing \texttt{InterruptedException} is considered a \textbf{bug}.
	\end{quote}
	
	\textbf{Deliverables:}
	\begin{itemize}[leftmargin=*]
		\item Thread dump before interruption
		\item Explanation of:
		\begin{itemize}[leftmargin=*]
			\item what interruption does \textit{not} fix (deadlock itself)
			\item why interruption must be handled explicitly
		\end{itemize}
		\item Updated scorecard with failure handling noted
	\end{itemize}
	
	\begin{quote}
		\textbf{Planted Insight (Intentional):}\\
		\textit{Failure isn't exceptional --- it's concurrent.}
	\end{quote}
	
	\subsubsection*{Deliverables (Mandatory visuals)}
	\begin{itemize}[leftmargin=*]
		\item \textbf{State diagram} of deadlock scenario
		\item Thread dump + interpretation
		\item Completed scorecard
	\end{itemize}
	
	% ------------------------------------------------------------
	\subsection*{Week 4 --- Memory Visibility \& Immutability}
	
	\subsubsection*{Concepts}
	\begin{itemize}[leftmargin=*]
		\item Visibility vs atomicity
		\item Caches, reordering, JMM
		\item \texttt{volatile}, happens-before
		\item Immutability as concurrency strategy
	\end{itemize}
	
	\subsubsection*{Lab 4.1--4.3}
	Broken counter / stale reads / DCL fix with \texttt{volatile}.
	
	\subsubsection*{Lab 4.4 (Concrete immutable snapshot context)}
	\textbf{Configuration cache updates without locks}
	\begin{itemize}[leftmargin=*]
		\item Shared mutable map causes visibility/race issues
		\item Fix using immutable snapshot + atomic swap of reference
	\end{itemize}
	
	\subsubsection*{Deliverables}
	\begin{itemize}[leftmargin=*]
		\item Short sequence diagram showing update/reader flow
		\item Completed scorecard
	\end{itemize}
	
	\subsubsection*{War Story Discussion (10 min)}
	Therac-25 (concurrency failure lessons).
	
	% ------------------------------------------------------------
	\subsection*{Week 5 --- \texttt{java.util.concurrent}: Safety, Liveness, Lock-Free}
	
	\subsubsection*{Concepts}
	\begin{itemize}[leftmargin=*]
		\item \texttt{ExecutorService}, thread pools
		\item Bounded queues + backpressure
		\item Atomics, CAS
		\item \texttt{ReentrantReadWriteLock} decision guidance:
		\begin{itemize}[leftmargin=*]
			\item Use when \textbf{read:write $>$ 10:1} and reads are non-trivial
		\end{itemize}
		\item \texttt{ThreadLocal} (useful vs abuse)
	\end{itemize}
	
	\subsubsection*{Lab 5.1}
	Bounded \texttt{BlockingQueue}: full vs empty behavior (backpressure).
	
	\subsubsection*{Lab 5.2}
	Thread pool sizing + bounded resource usage.
	
	\subsubsection*{Lab 5.3 (Lock-free programming required)}
	\textbf{CAS counter vs \texttt{synchronized} counter}
	\begin{itemize}[leftmargin=*]
		\item Compare throughput and correctness under contention
		\item Discuss why atomics exist
	\end{itemize}
	
	\subsubsection*{Lab 5.4}
	Read-heavy workload: \texttt{ReentrantReadWriteLock} vs \texttt{synchronized}.
	
	\subsubsection*{Lab 5.5 (ThreadLocal required)}
	\textbf{Request context propagation}
	\begin{itemize}[leftmargin=*]
		\item Attach userID/requestID via \texttt{ThreadLocal}
		\item Show correct use + ``leak'' hazard (not clearing)
	\end{itemize}
	
	\subsubsection*{Deliverables}
	\begin{itemize}[leftmargin=*]
		\item Completed scorecard + liveness note (starvation risks)
		\item Measurement table + tradeoff paragraph
	\end{itemize}
	
	\subsubsection*{Pitfalls Cheat Sheet (Release 2)}
	\begin{itemize}[leftmargin=*]
		\item Don't oversize pools
		\item Don't use unbounded queues in production
	\end{itemize}
	
	% ------------------------------------------------------------
	\subsection*{Week 6 --- Parallelism \& Fork/Join (Granularity + Break-Even)}
	
	\subsubsection*{Concepts}
	\begin{itemize}[leftmargin=*]
		\item Task decomposition, work stealing
		\item ForkJoinPool subtleties
		\item \textbf{Pitfall:} blocking I/O inside ForkJoinPool starves workers
	\end{itemize}
	
	\subsubsection*{Lab 6.1}
	Fork/Join vs sequential (compute task).
	
	\subsubsection*{Lab 6.2 (Numeric target)}
	\textbf{Determine minimum array size where Fork/Join beats sequential}
	\begin{itemize}[leftmargin=*]
		\item Plot/record break-even point
		\item Explain coordination overhead
	\end{itemize}
	
	\subsubsection*{Lab 6.3 (Parallel streams anti-example + why)}
	Show slowdown and explain:
	\begin{itemize}[leftmargin=*]
		\item boxing overhead
		\item lambda allocation
		\item spliterator/partitioning inefficiency
		\item thread management overhead
	\end{itemize}
	
	\subsubsection*{Case Study (Revisit Redis)}
	Redis v6+ I/O threads: where threads actually appear and why.
	
	\subsubsection*{Deliverables}
	\begin{itemize}[leftmargin=*]
		\item Completed scorecard
		\item Break-even calculation and explanation
	\end{itemize}
	
	% ------------------------------------------------------------
	\subsection*{Week 7 --- Async Design, Testing, and Anti-Patterns}
	
	\subsubsection*{Concepts}
	\begin{itemize}[leftmargin=*]
		\item \texttt{CompletableFuture} pipelines
		\item Timeouts/fallbacks
		\item Testing concurrent code (full treatment):
		\begin{itemize}[leftmargin=*]
			\item Stress testing: \textbf{JCStress overview and usage}
			\item Deterministic replay techniques (conceptual)
			\item Mock time/delays to reproduce races
			\item Reproduction strategy: \texttt{Thread.yield()}, tight loops, randomized scheduling
		\end{itemize}
	\end{itemize}
	
	\subsubsection*{Lab 7.1}
	Async service simulation + timeout/fallback.
	
	\subsubsection*{Lab 7.2}
	Testing trap: passes 1000, fails at 1001.  
	\textbf{Add reproduction strategy:} yield + iteration amplification.
	
	\subsubsection*{Lab 7.3}
	JCStress-style stress test (guided).
	
	\subsubsection*{Lab 7.4 --- Fan-Out / Fan-In Async Pipeline (Required)}
	
	\textbf{Objective:} teach structured async composition and aggregation --- not ``fire-and-forget'' chaos.
	
	\textbf{Scenario:} a request is split into multiple independent async sub-tasks (fan-out), whose results must be combined into a single response (fan-in).
	
	\textbf{Tasks:}
	\begin{itemize}[leftmargin=*]
		\item Implement fan-out using \texttt{CompletableFuture.supplyAsync}
		\item Combine results using:
		\begin{itemize}[leftmargin=*]
			\item \texttt{allOf()} for synchronization
			\item Explicit result aggregation
		\end{itemize}
		\item Inject:
		\begin{itemize}[leftmargin=*]
			\item One slow task
			\item One failing task
		\end{itemize}
		\item Add timeout and fallback for the whole pipeline
	\end{itemize}
	
	\textbf{Deliverables:}
	\begin{itemize}[leftmargin=*]
		\item Sequence diagram of fan-out / fan-in flow
		\item Explanation of:
		\begin{itemize}[leftmargin=*]
			\item how completion is coordinated
			\item how failure propagates
		\end{itemize}
		\item Completed scorecard
	\end{itemize}
	
	\begin{quote}
		\textbf{Core realization:}\\
		\textit{Async systems are pipelines, not threads with callbacks.}
	\end{quote}
	
	\subsubsection*{Lab 7.5 (Required: NIO counterpoint)}
	\textbf{Single-threaded async I/O (Java NIO) vs multi-threaded blocking}
	\begin{itemize}[leftmargin=*]
		\item Same hardware
		\item Measure throughput/latency
		\item Prove: async $\neq$ more threads
	\end{itemize}
	
	\subsubsection*{Anti-Patterns Module (with mini-examples)}
	\begin{itemize}[leftmargin=*]
		\item Thread-per-request
		\item Unbounded queues
		\item Missing timeouts
		\item Silent failure swallowing
	\end{itemize}
	
	\subsubsection*{Deliverables}
	\begin{itemize}[leftmargin=*]
		\item Completed scorecard
		\item Sequence diagram of async pipeline
	\end{itemize}
	
	% ------------------------------------------------------------
	\subsection*{Week 8 --- Midterm: Debugging + Kafka Mental Model Preview}
	
	\subsubsection*{Midterm Submission (Tool usage required)}
	\begin{itemize}[leftmargin=*]
		\item Thread dump
		\item Diagnosis report
		\item Fix
		\item Performance tradeoff analysis
		\item 5-minute recorded explanation to technical audience
	\end{itemize}
	
	\textbf{Rubric:} 1) real bug, 2) mechanism, 3) cost.
	
	\subsubsection*{Kafka Preview (20-minute lecture, no lab)}
	\textbf{``Logs are the universal abstraction.''}
	\begin{itemize}[leftmargin=*]
		\item Log as history, replay, recovery, audit, decoupling
		\item Why distributed logs show up everywhere (Kafka, event sourcing)
	\end{itemize}
	
	\subsubsection*{Pitfalls Cheat Sheet (Release 3)}
	\begin{itemize}[leftmargin=*]
		\item Measure before optimizing
		\item Async $\neq$ faster
		\item Fix mechanism, not symptom
	\end{itemize}
	
	% ============================================================
	\section*{Phase 2 --- Language Models \& Limits (Week 9)}
	
	\subsection*{Week 9 --- Python GIL + Brief Survey of Alternatives}
	
	\subsubsection*{Concepts}
	\begin{itemize}[leftmargin=*]
		\item GIL reality: CPU-bound threads don't scale
		\item Multiprocessing tradeoffs
		\item 15-minute contrast:
		\begin{itemize}[leftmargin=*]
			\item Go goroutines (scheduler + channels)
			\item Erlang actors (message passing model)
		\end{itemize}
	\end{itemize}
	
	\subsubsection*{Lab 9.1}
	Python threads vs multiprocessing (CPU-bound).
	
	\subsubsection*{Lab 9.2 (Visual proof)}
	\texttt{htop} / \texttt{py-spy}:
	\begin{itemize}[leftmargin=*]
		\item 4 threads $\rightarrow$ 1 core active in Python
		\item 4 Java threads $\rightarrow$ multiple cores
	\end{itemize}
	
	\subsubsection*{War Story Discussion (10 min)}
	Knight Capital (timing + deployment + race-type failures).
	
	% ============================================================
	\section*{Phase 3 --- Distributed Systems (Weeks 10--14)}
	
	% ------------------------------------------------------------
	\subsection*{Week 10 --- Network = Concurrency + Failure (Phase 3 Start)}
	
	\subsubsection*{Concepts}
	\begin{itemize}[leftmargin=*]
		\item Networks are slow queues with partial failure
		\item No shared memory
		\item Retry amplifies load
	\end{itemize}
	
	\subsubsection*{Three Laws of Networks}
	\begin{enumerate}[leftmargin=*]
		\item Latency is never zero
		\item Bandwidth is never infinite
		\item Failures are partial and invisible
	\end{enumerate}
	
	\subsubsection*{Tool Introduction (Required): Toxiproxy}
	Used in Weeks 10--13:
	\begin{itemize}[leftmargin=*]
		\item latency injection
		\item bandwidth limit
		\item connection reset/timeouts
	\end{itemize}
	
	\subsubsection*{Lab 10.1}
	Local \texttt{BlockingQueue} $\rightarrow$ socket queue (same business logic).  
	Then use Toxiproxy to inject:
	\begin{itemize}[leftmargin=*]
		\item latency
		\item packet loss/reset
	\end{itemize}
	
	\textbf{Deliverables:}
	\begin{itemize}[leftmargin=*]
		\item Completed scorecard
		\item Measurement summary
	\end{itemize}
	
	% ------------------------------------------------------------
	\subsection*{Week 11 --- Client--Server, Serialization, Versioning + Peer Review}
	
	\subsubsection*{Concepts}
	\begin{itemize}[leftmargin=*]
		\item TCP sockets
		\item serialization boundary = no shared objects
		\item defensive copying, immutability
		\item versioning failures
	\end{itemize}
	
	\subsubsection*{Lab 11.1}
	Multi-client server, immutability fix.
	
	\subsubsection*{Lab 11.2 (Versioning: specify how)}
	Choose one:
	\begin{itemize}[leftmargin=*]
		\item Java serialization with \texttt{serialVersionUID} and compatibility handling, OR
		\item Protocol Buffers with field evolution (recommended)
	\end{itemize}
	
	\subsubsection*{Peer Review Checkpoint (Required)}
	Students swap Week 10 socket code:
	\begin{itemize}[leftmargin=*]
		\item Find concurrency bug or failure scenario
		\item Provide written review (what breaks, evidence, suggested fix)
	\end{itemize}
	
	\subsubsection*{Deliverables}
	\begin{itemize}[leftmargin=*]
		\item Completed scorecard
		\item \textbf{Sequence diagram} showing message flow
	\end{itemize}
	
	\subsubsection*{Pitfalls Cheat Sheet (Release 4)}
	\begin{itemize}[leftmargin=*]
		\item Don't assume network reliability
		\item Always time out network calls
	\end{itemize}
	
	% ------------------------------------------------------------
	\subsection*{Week 12 --- Kafka: Partitions, Consumer Groups, Backpressure}
	
	\subsubsection*{Concepts}
	\begin{itemize}[leftmargin=*]
		\item Kafka as distributed log (not ``Kafka course'')
		\item Partitions: ordering vs parallelism
		\item Consumer groups + rebalancing
		\item Backpressure and lag
	\end{itemize}
	
	\subsubsection*{Lab 12.1}
	Partitions + same key routing + ordering effects.
	
	\subsubsection*{Lab 12.2 (Consumer groups)}
	2 consumers, 4 partitions:
	\begin{itemize}[leftmargin=*]
		\item observe assignment
		\item trigger rebalance and observe behavior
	\end{itemize}
	
	\subsubsection*{Lab 12.3 (Backpressure)}
	Flood topic with 10k msg/s, 1 slow consumer:
	\begin{itemize}[leftmargin=*]
		\item measure lag growth
		\item discuss mitigation (more partitions, faster consumers, batch, retention)
	\end{itemize}
	
	\subsubsection*{Partition Decision Exercise (Constrained)}
	Given:
	\begin{itemize}[leftmargin=*]
		\item 1M msg/s
		\item 3 consumers
		\item 100ms max lag target
	\end{itemize}
	Students must propose partition count with calculations and assumptions.
	
	\subsubsection*{Case Study}
	Kafka at LinkedIn (design motivations).
	
	% ------------------------------------------------------------
	\subsection*{Week 13 --- RPC, Cascading Failure, Tracing, Idempotency, Clock Drift}
	
	\subsubsection*{Concepts}
	\begin{itemize}[leftmargin=*]
		\item gRPC: contracts and timeouts
		\item retries and their dangers
		\item \textbf{idempotency:} retries require idempotent operations
		\item distributed tracing without fancy tools: request IDs + log correlation
		\item log aggregation basics: grep across logs + timestamp correlation
		\item clock drift: time is the enemy
	\end{itemize}
	
	\subsubsection*{Circuit Breaker (No instructor dependency)}
	Students implement using one of:
	\begin{itemize}[leftmargin=*]
		\item \textbf{Resilience4j} (recommended), OR
		\item Pseudocode spec (state machine: closed/open/half-open)
	\end{itemize}
	
	\subsubsection*{Lab 13.1 --- Cascading Failure (Explicit)}
	\begin{enumerate}[leftmargin=*]
		\item Healthy baseline (measure)
		\item Inject 10\% timeouts via Toxiproxy
		\item Observe thread pool exhaustion
		\item Trace failure propagation
		\item Add circuit breaker
		\item Add idempotency fix for retry-safe operations
	\end{enumerate}
	
	\subsubsection*{Lab 13.2 --- Tracing Reconstruction}
	Log trace through 3-service chain:
	\begin{itemize}[leftmargin=*]
		\item propagate request IDs
		\item reconstruct timeline from logs (manual trace analysis)
	\end{itemize}
	
	\subsubsection*{Lab 13.3 --- Clock Drift Demo}
	Skew clocks; show timestamp ordering violation and why correlation must be careful.
	
	% ------------------------------------------------------------
	\subsection*{Week 14 --- Final Synthesis, Presentations, Tool Evaluation}
	
	\subsubsection*{Portfolio (80\%) --- Checkpoint Weights}
	\begin{itemize}[leftmargin=*]
		\item Week 3: \textbf{10\%}
		\item Week 5: \textbf{15\%}
		\item Week 8 (Midterm): \textbf{25\%}
		\item Week 14 (Final): \textbf{30\%}
	\end{itemize}
	
	Final must include:
	\begin{itemize}[leftmargin=*]
		\item explicit failure guarantees (timeouts, retry behavior)
		\item tracing evidence
		\item idempotency handling
		\item completed scorecard
	\end{itemize}
	
	\subsubsection*{Architecture Decision Memo (20\%) --- Rubric}
	\begin{itemize}[leftmargin=*]
		\item Correctness of tradeoff analysis: \textbf{40\%}
		\item Justification of choice: \textbf{30\%}
		\item Consideration of failure modes: \textbf{20\%}
		\item Clarity: \textbf{10\%}
	\end{itemize}
	
	Memo uses decision matrix:
	\begin{itemize}[leftmargin=*]
		\item latency requirements
		\item ordering guarantees
		\item failure modes
		\item operational complexity
	\end{itemize}
	
	\subsubsection*{Presentation (Week 14)}
	\textbf{5 minutes} to a technical audience:
	\begin{itemize}[leftmargin=*]
		\item explain one failure from your distributed system
		\item show evidence (logs/traces/metrics)
		\item explain fix and tradeoff
	\end{itemize}
	
	\subsubsection*{Further Study (Specified)}
	Annotated bibliography (3--5 items each):
	\begin{itemize}[leftmargin=*]
		\item Actor models
		\item CRDTs
		\item Consensus (Raft/Paxos)
		\item Advanced tracing/observability
		\item Formal concurrency testing
	\end{itemize}
	
	\subsubsection*{Tool Evaluation Framework (Operationalizing ``tools change'')}
	Checklist for evaluating any new concurrency tool/framework:
	\begin{itemize}[leftmargin=*]
		\item What guarantees does it provide?
		\item What failure modes does it hide or introduce?
		\item What is the mental model?
		\item How does it behave under overload?
		\item How do you debug it?
	\end{itemize}
	
	\subsubsection*{War Story Discussion (10 min)}
	AWS S3 outage synthesis + ``what would you instrument?''
	
	% ============================================================
	\section*{Case Studies + Discussion Component (Required)}
	\begin{itemize}[leftmargin=*]
		\item Phase 1: Therac-25 (Week 4 discussion)
		\item Phase 2: Knight Capital (Week 9 discussion)
		\item Phase 3: AWS S3 outage (Week 14 discussion)
	\end{itemize}
	
	% ============================================================
	\section*{Progressive Visual Requirements}
	\begin{itemize}[leftmargin=*]
		\item Week 3: \textbf{state diagram} (deadlock)
		\item Week 8: \textbf{sequence diagram} (async pipeline)
		\item Week 14: \textbf{system sequence + failure propagation diagram}
	\end{itemize}
	
	% ============================================================
	\section*{Final Statement}
	
	This course does not train you to ``use concurrency libraries.''  
	It trains you to \textbf{diagnose broken systems}.
	
	When your future production system freezes, corrupts data, or collapses under load, you will not panic.  
	You will take a thread dump, trace the failure, identify the mechanism, and fix it professionally.
	
\end{document}
